\section{Introduction}
\label{sec:introduction}

Time-series forecasting has converged on a foundation-model recipe
inherited from language modelling: patched inputs, a transformer
backbone, and a head that forecasts the next patch in value space.
TimesFM~\citep{das2024timesfm} is the canonical example, and recent
successors~\citep{ansari2024chronos, woo2024moirai, shi2024timemoe,
liu2025sundial, liu2026timers1} extend the recipe along various axes
while keeping the value-space target. This choice is structurally
suboptimal: the observed series is one realisation of a stochastic
process, so a value-space loss asks the model to reproduce the
intrinsic noise of that process. No model can succeed at that task,
and the training loss never reaches zero, even on processes whose
predictable component the model captures perfectly.

A second line of work shifts the prediction target from raw values to
a learned representation. The Joint Embedding Predictive
Architecture (JEPA)~\citep{assran2023ijepa}, originally introduced
for vision, encodes a context block with one network and a target
block with another, and trains a predictor to map the context latent
to the target latent. The recent time-series
adaptations~\citep{verdenius2024latpfn, ennadir2025tsjepa} apply the
same idea to univariate series. By predicting an abstract feature of
the future rather than its raw values, the model is no longer
punished for missing the noise. The siamese architecture used by
JEPA, however, is at risk of collapse by construction: the constant
function is a global minimum of its loss. Contrastive learning with
explicit negatives is the classical anti-collapse remedy, and has
been used at scale for time-series representation learning
(CPC~\citep{oord2018cpc}, TS2Vec~\citep{yue2022ts2vec}), but to our
knowledge no published time-series foundation model uses it as a
primary forecasting objective.

We ask whether a time-series forecasting model can be trained
directly in latent space with an explicit contrastive objective, in
such a way that the latent it learns encodes the underlying
data-generating process. The rest of the paper proposes a candidate
architecture and tests this question on a controlled synthetic
distribution, where the generating parameters are known and
recoverability can be checked directly.

The contributions of the paper are the following.
\begin{itemize}
\item A training framework, \emph{Contrastive Forecasting}, in which
  an encoder and a causal forecaster are trained jointly with an
  InfoNCE objective on cross-temporal, cross-channel and cross-batch
  negatives (Section~\ref{sec:method}). The framework is decoupled
  from the choice of encoder and forecaster architecture.
\item Empirical evidence that the resulting latents encode the
  parameters of the data-generating process. On synthetic ARMA, a
  frozen-backbone recovery head retrieves the AR and MA coefficients
  with an improvement ratio well above the all-zero baseline
  (Sections~\ref{sec:experiments} and~\ref{sec:results}).
\item An architecture-search trail for both the encoder and the
  forecaster, with the design decisions reported in
  Appendix~\ref{sec:appendix-backbone-search}.
\end{itemize}
